\phantomsection
\subsection*{M4F. Experimental Mead}
\addcontentsline{toc}{subsection}{M4F. Experimental Mead}

\textit{Um Experimental Mead é um hidromel que não se enquadra em nenhuma outra categoria ou estilo de hidromel. Nenhum hidromel pode ser considerado ``fora de estilo'' nesta categoria, a menos que se enquadre em outro estilo de hidromel existente.}

\textbf{Impressão Geral}: Depende da natureza experimental declarada, mas deve ser reconhecível como hidromel, apresentar componentes experimentais identificáveis e ser agradável de beber.

\textbf{Aroma, Aparência, Sabor, Sensação na Boca}: Dependem inteiramente da natureza experimental do hidromel e dos ingredientes especiais, processos ou equilíbrio declarados. Tanto o caráter especial quanto o hidromel-base devem ser identificáveis e complementar um ao outro. O equilíbrio geral e a facilidade de beber do produto final são os principais fatores de sucesso deste estilo. A criatividade pode ser valorizada, desde que o hidromel continue agradável de beber.

\textbf{Ingredientes}: Dependem da natureza experimental.

\textbf{Comentários}: Alguns exemplos desta categoria incluem:

\begin{itemize}

\item Hidromel que combina vários estilos de hidromel, exceto quando a combinação se enquadrar em outro estilo, como M3C Fruit and Spice Mead.
\item Hidromel com fontes adicionais de açúcares fermentáveis ou não fermentáveis (por exemplo, xarope de bordo, melaço, açúcar mascavo, néctar de agave, lactose).
\item Hidromel com ingredientes adicionais (por exemplo, licores, destilados, defumação, produtos derivados de leite, lúpulo).
\item Hidromel produzido com processos, métodos ou técnicas alternativas (por exemplo, congelamento, defumação, acidificação).
\item Hidromel fermentado com leveduras ou bactérias não tradicionais que produzam um efeito perceptível (por exemplo, Brettanomyces, Lactobacillus, lambic belga, sake).
\item Hidromel regional, tradicional, histórico ou nativo que não esteja contemplado nestas diretrizes.
\item Hidromel ``conceitual'', elaborado para remeter a um coquetel, sobremesa ou outro alimento ou bebida, exceto quando se enquadrar em um estilo já estabelecido.
\item Hidromel que, nos demais aspectos, atenda às descrições de um estilo existente nas diretrizes, mas que esteja perceptivelmente fora dos parâmetros estabelecidos para esse estilo ou utilize ingredientes não tradicionais.

\end{itemize}

Esta lista é apenas representativa e não deve ser interpretada como uma limitação dos Experimental Meads a esses exemplos. A natureza experimental não deve incluir nada que seja ilegal na jurisdição onde o hidromel é produzido ou julgado, nem qualquer elemento que possa representar risco aos juízes.

\textbf{Instruções para Inscrição}: Os participantes devem especificar os níveis de dulçor, carbonatação e teor alcoólico. O participante deve descrever a natureza experimental do hidromel com detalhes suficientes para justificar sua inscrição neste estilo.

\textbf{Exemplos Comerciais}: Lost Cause Goblin King, Lost Cause Lemon's Cello, Prairie Rose Mojito, Redstone Nectar of the Hops, Samuel Adams Honey Queen, St. Ambrose Razzputin, White Winter Black Harbor